\section{文档目的与范围}

本文规定「寻迹」校园失物招领平台的配置管理、版本控制、持续集成、部署和运维方法。
文档的目标是让开发成员使用一致的配置名称和版本规则，让部署人员能够按固定步骤启动
系统，并让运行期间的检查、备份和变更都有明确流程。

\begin{table}[H]
  \centering
  \caption{配置与运维管理范围}
  \begin{tabularx}{\textwidth}{L{3.0cm}L{4.1cm}Y}
    \toprule
    管理方向 & 采用的机制 & 管理目的 \\
    \midrule
    配置管理 & 环境变量、示例模板、Pydantic Settings、Vite 构建参数 & 分离代码、环境差异和敏感值 \\
    版本控制 & Git 分支、规范化提交、版本标签、Alembic revision & 追踪应用与数据库的每次变化 \\
    持续集成 & GitHub Actions 并行执行静态检查、测试和构建 & 在合并前使用统一规则检查各模块 \\
    部署管理 & Dockerfile、Compose、健康检查和部署清单 & 以相同步骤构建和启动六个服务 \\
    运行维护 & 日志、运行指标、巡检、备份、恢复和回滚流程 & 保持服务与数据处于可管理状态 \\
    \bottomrule
  \end{tabularx}
\end{table}

\section{配置管理}

\subsection{配置分层}

项目遵循「代码与配置分离、敏感值与模板分离、不同环境使用相同变量名」的原则。配置按
以下四层组织：

\begin{enumerate}
  \item \textbf{代码默认值：}只保存通用且非敏感的运行参数，例如 API 前缀、默认页大小和
  本地端口。
  \item \textbf{示例模板：}列出部署所需变量、说明和示例格式，空出密码、token 和私钥。
  \item \textbf{环境配置：}部署者从模板生成本地环境文件，为当前环境填写真实值；环境
  文件不进入版本库，也不复制到镜像层。
  \item \textbf{启动注入：}Compose 将环境变量传入容器，Backend 和 AI Service 通过
  Pydantic Settings 读取；前端 API 基址通过 Vite 构建参数写入静态产物。
\end{enumerate}

进程环境变量优先于环境文件中的同名值，便于部署平台在不改动仓库文件的情况下覆盖配置。
配置名称保持大写下划线形式，布尔值使用 \term{true/false}，列表值采用英文逗号分隔，
时间参数在名称中明确使用秒、分钟或小时。

\subsection{配置项分类}

\small
\begin{longtable}{L{2.6cm}L{5.2cm}L{6.5cm}}
  \caption{主要配置类别}\label{tab:configuration-categories}\\
  \toprule
  类别 & 代表变量 & 管理规则 \\
  \midrule
  \endfirsthead
  \multicolumn{3}{c}{\tablename\ \thetable（续）}\\
  \toprule
  类别 & 代表变量 & 管理规则 \\
  \midrule
  \endhead
  \bottomrule
  \endlastfoot
  数据库 & \texttt{DB\_HOST}、\texttt{DB\_PORT}、\texttt{DB\_NAME}、\texttt{DB\_USER}、\texttt{DB\_PASSWORD} & 库名和账号按环境区分，密码由部署环境注入 \\
  身份鉴权 & \texttt{JWT\_SECRET\_KEY}、\texttt{JWT\_EXPIRE\_DAYS} & 共享环境使用独立随机密钥，不沿用示例值 \\
  管理员初始化 & \texttt{BOOTSTRAP\_ADMIN\_ENABLED}、管理员账号、密码和手机号 & 仅首次初始化时开启，账号建立后关闭开关 \\
  对象存储 & MinIO endpoint、access key、secret key、bucket、签名时长 & bucket 保持私有，浏览器和 AI 使用不同有效期的签名 \\
  AI 服务 & \texttt{AI\_SERVICE\_BASE\_URL}、\texttt{AI\_SERVICE\_TOKEN}、timeout & 主后端和 AI 服务配置同一随机 token，服务间调用携带鉴权头 \\
  模型供应商 & API key、base URL、文本模型、向量模型和视觉模型 & key 由部署环境保管，模型名按环境切换 \\
  Web 访问 & \texttt{CORS\_ALLOWED\_ORIGINS}、\texttt{VITE\_API\_BASE\_URL} & 后端使用显式来源列表，前端构建参数与部署域名一致 \\
  后台任务 & 匹配间隔、任务轮询、最大尝试次数、退避和锁超时 & 时间单位写入变量名，修改时同步说明和运维记录 \\
  本地演示 & \texttt{DEBUG}、\texttt{SMS\_DEBUG\_ENABLED}、\texttt{SMS\_DEMO\_PHONES} & 调试验证码仅对明确手机号列表生效，共享环境关闭调试开关 \\
\end{longtable}
\normalsize

\subsection{启动校验}

Compose 对 MySQL root 密码、业务数据库账号、JWT 密钥、管理员账号和 MinIO 凭据使用
必填变量表达式，缺少关键值时停止配置展开。Backend 在非 Debug 环境拒绝示例 JWT 密钥
和示例管理员密码，并拒绝通配 CORS；AI Service 要求服务 token 至少 32 个字符，只有
显式本地模式允许绕过。启动校验把配置规则放到程序入口，避免依赖部署者记忆。

\subsection{环境划分}

\begin{table}[H]
  \centering
  \caption{环境配置策略}
  \begin{tabularx}{\textwidth}{L{2.7cm}L{3.5cm}Y Y}
    \toprule
    环境 & 用途 & 配置特点 & 数据策略 \\
    \midrule
    本地开发 & 单人开发与调试 & 本机端口、可开启精确短信白名单、独立密钥 & 每名成员使用独立测试数据 \\
    团队联调 & 跨模块接口联调 & 固定 API 地址、统一 CORS、关闭调试返回值 & 使用团队测试账号与可重建数据 \\
    课程演示 & 答辩和功能展示 & 冻结版本、固定端口、独立随机密钥 & 预置演示数据，部署前制作快照 \\
    正式部署 & 长期提供服务 & HTTPS、最小端口暴露、独立服务账号 & 执行周期备份和恢复校验 \\
    \bottomrule
  \end{tabularx}
\end{table}

环境之间复用变量名称和部署结构，但不复用数据库、对象 bucket 或访问凭据。配置文件只
保存当前环境需要的值；从联调环境切换到演示环境时，通过替换环境文件和镜像版本完成，
不在源代码中加入环境判断。

\subsection{配置变更流程}

\begin{enumerate}
  \item 提交者说明变量名称、类型、默认行为、取值范围和适用环境。
  \item 同步更新 Settings、Compose 映射和示例模板；前端变量同时更新构建说明。
  \item 本地执行配置展开与应用启动检查，确认缺省值和必填值行为。
  \item 通过 Pull Request 评审变量命名、敏感性和兼容性。
  \item 部署时先在联调环境应用新配置，再按相同名称更新演示或正式环境。
  \item 记录变更时间、应用版本、变量名称和操作者，不记录秘密值本身。
\end{enumerate}

密码、JWT 密钥、AI token、数据库口令和对象存储 secret 不写入 issue、提交信息、日志或
截图。需要更换凭据时，先生成新值并更新服务端，再更新调用方，确认新值生效后撤销旧值。

\section{版本控制}

\subsection{Git 分支与提交}

仓库以 \term{main} 保存稳定交付版本，以 \term{develop} 汇集团队增量；功能和修复分别
使用 \term{feature/<name>}、\term{fix/<name>}。每个 Pull Request 聚焦一个功能或一个问题，
并关联需求编号或任务。提交信息采用简短的 Conventional Commit，例如：

\begin{verbatim}
feat(FR-02): add lost item publish api
fix(FR-05): correct claim review transition
test(match): cover fallback calculation
docs: update deployment instructions
chore(ci): adjust compose check
\end{verbatim}

\begin{table}[H]
  \centering
  \caption{版本控制对象与标识}
  \begin{tabularx}{\textwidth}{L{3.2cm}L{4.0cm}Y}
    \toprule
    对象 & 版本标识 & 管理方法 \\
    \midrule
    源代码与文档 & Git commit SHA & 同一提交保存相关代码、迁移、配置模板和说明 \\
    稳定版本 & Git tag & 标签指向经评审的固定提交，不移动已有标签 \\
    Python 依赖 & uv.lock & Backend 与 AI Service 分别冻结解析版本 \\
    前端依赖 & pnpm-lock.yaml & 用户端和管理端分别执行冻结安装 \\
    数据库结构 & Alembic revision & 每次结构变化形成有顺序的迁移版本 \\
    容器镜像 & 应用名加 commit SHA 或版本号 & 部署记录能够反查源码与依赖 \\
    \bottomrule
  \end{tabularx}
\end{table}

\path{.gitattributes} 统一普通文本为 LF，Windows 命令脚本使用 CRLF，二进制文件不参与
行尾转换。\path{.gitignore} 排除环境文件、依赖目录、缓存、日志和构建中间产物，最终
课程 PDF 作为交付文件单独纳入版本控制。

\subsection{数据库版本}

数据库使用 SQLAlchemy 模型描述当前结构，使用 Alembic 保存从旧版本到新版本的变化。
迁移文件按单一 revision 链排列，Backend 容器在启动 Uvicorn 前运行
\term{alembic upgrade head}。手写 Schema 文件不作为新环境初始化入口，避免 ORM 与独立
建表脚本产生两套结构定义。

\begin{longtable}{L{3.0cm}L{7.1cm}L{4.3cm}}
  \caption{数据库变更管理规则}\label{tab:database-version-rules}\\
  \toprule
  变更类型 & 迁移要求 & 评审内容 \\
  \midrule
  \endfirsthead
  \multicolumn{3}{c}{\tablename\ \thetable（续）}\\
  \toprule
  变更类型 & 迁移要求 & 评审内容 \\
  \midrule
  \endhead
  \bottomrule
  \endlastfoot
  新增表或字段 & 明确类型、nullability、默认值和索引 & ORM 与数据库设计是否一致 \\
  新增唯一约束 & 先整理历史重复记录，再建立约束 & 数据转换顺序和业务唯一性 \\
  修改状态值 & 给出旧值到新值的映射 & Service 状态机与枚举是否同步 \\
  删除字段 & 先停止代码读写，再迁移数据并删除 & 新旧应用是否存在兼容窗口 \\
  历史数据库接入 & 备份并逐项比较表、列、索引和约束 & 结构等价后才允许记录当前 revision \\
\end{longtable}

迁移提交同时修改 ORM、数据库说明和相关测试。自动生成的迁移只作为初稿，提交者需要
检查数据语句和 downgrade 行为。部署记录保存变更前后 revision，便于确认环境所处版本。

\section{持续集成}

\subsection{触发与并行任务}

GitHub Actions 在 \term{main}、\term{develop} 的 push 和 Pull Request 上运行。相同分支
出现新提交时，旧运行被取消；Backend、AI Service、User App、Admin Web 和 Docker Build
五类任务并行执行。所有任务采用仓库锁文件安装依赖，因而开发者可以在本地复现同一组
命令。

\begin{longtable}{L{2.9cm}L{7.2cm}L{4.3cm}}
  \caption{持续集成任务}\label{tab:configuration-ci-jobs}\\
  \toprule
  任务 & 执行步骤 & 管理作用 \\
  \midrule
  \endfirsthead
  \multicolumn{3}{c}{\tablename\ \thetable（续）}\\
  \toprule
  任务 & 执行步骤 & 管理作用 \\
  \midrule
  \endhead
  \bottomrule
  \endlastfoot
  Backend & uv 安装；Ruff lint/format；Mypy；Pytest 与覆盖率 & 检查主后端规范、类型和业务逻辑 \\
  AI Service & uv 安装；Ruff lint/format；Mypy；Pytest & 检查独立服务及其 HTTP 契约 \\
  User App & pnpm 冻结安装；vue-tsc；Vitest；Vite build & 检查用户端类型、逻辑和生产构建 \\
  Admin Web & pnpm 冻结安装；vue-tsc；Vitest；Vite build & 检查管理端类型、逻辑和生产构建 \\
  Docker Build & Compose config；Backend 与 AI 镜像构建 & 检查编排配置和镜像构建入口 \\
\end{longtable}

流水线失败时，提交者根据任务日志在原分支修正，并由新提交触发重新检查。覆盖率 XML
作为 Backend 任务产物保存；Compose 配置检查使用专用占位值，只验证环境变量映射和
编排语法，不读取部署环境中的真实凭据。

\subsection{版本进入部署的条件}

进入演示或正式部署的版本应指向固定 commit 或 tag，并满足：Pull Request 已评审；持续
集成任务成功；环境模板与新增变量同步；数据库迁移随版本提交；部署者能够从锁文件构建
镜像。部署记录使用 commit SHA、镜像标识、Alembic revision 和配置版本共同描述一次
发布，避免只记录「最新版」。

\section{部署架构}

\subsection{六服务拓扑}

\begin{figure}[H]
  \centering
  \begin{tikzpicture}[node distance=10mm and 13mm]
    \node[actor,minimum width=4.0cm] (user) {普通用户浏览器};
    \node[actor,minimum width=4.0cm,right=36mm of user] (admin) {管理员浏览器};
    \node[component,below=of user] (userweb) {user-web\\Nginx :18080};
    \node[component,below=of admin] (adminweb) {admin-web\\Nginx :18081};
    \coordinate (webmid) at ($(userweb)!0.5!(adminweb)$);
    \node[component,below=14mm of webmid,minimum width=4.2cm] (backend) {FastAPI Backend\\Uvicorn :8080};
    \node[datastore,below left=16mm and 27mm of backend] (mysql) {MySQL 8.4\\:3306};
    \node[datastore,below=16mm of backend] (minio) {MinIO\\:9000/:9001};
    \node[component,below right=16mm and 27mm of backend] (ai) {AI Service\\:5000};
    \node[note,below=10mm of mysql] (mysqlvol) {mysql-data\\持久卷};
    \node[note,below=10mm of minio] (miniovol) {minio-data\\持久卷};
    \node[note,below=10mm of ai] (provider) {可选模型服务};
    \draw[flow] (user) -- (userweb);
    \draw[flow] (admin) -- (adminweb);
    \draw[flow] (userweb) -- node[left,font=\zihao{-5}]{/api/v1} (backend);
    \draw[flow] (adminweb) -- node[right,font=\zihao{-5}]{/api/v1} (backend);
    \draw[flow] (backend) -- node[above left,font=\zihao{-5}]{async SQL} (mysql);
    \draw[flow] (backend) -- node[right,font=\zihao{-5}]{私有对象/签名} (minio);
    \draw[flow] (backend) -- node[above right,font=\zihao{-5}]{服务 token} (ai);
    \draw[dashedflow] (ai) -- (minio);
    \draw[dashedflow] (ai) -- (provider);
    \draw[flow] (mysql) -- (mysqlvol);
    \draw[flow] (minio) -- (miniovol);
  \end{tikzpicture}
  \caption{完整栈部署拓扑}
  \label{fig:deployment-topology}
\end{figure}

两个前端由 Nginx 提供静态文件，并将 \path{/api/v1} 转发到 Backend。Backend 管理业务、
数据库迁移和对象签名；AI Service 作为独立 HTTP 服务运行，通过 service token 鉴权；
MySQL 和 MinIO 使用命名卷保存数据。数据库、对象存储和 AI 端口在本地配置中绑定回环
地址，外部用户只访问前端和 Backend 入口。

\subsection{服务清单}

\small
\begin{longtable}{L{2.0cm}L{2.9cm}L{2.4cm}L{3.3cm}L{3.0cm}}
  \caption{Compose 服务与启动关系}\label{tab:compose-services}\\
  \toprule
  服务 & 镜像或构建 & 默认端口 & 数据与健康检查 & 启动关系 \\
  \midrule
  \endfirsthead
  \multicolumn{5}{c}{\tablename\ \thetable（续）}\\
  \toprule
  服务 & 镜像或构建 & 默认端口 & 数据与健康检查 & 启动关系 \\
  \midrule
  \endhead
  \bottomrule
  \endlastfoot
  mysql & MySQL 8.4 & 127.0.0.1:3306 & mysql-data；mysqladmin ping & 先启动并达到 healthy \\
  minio & 固定日期镜像 & 127.0.0.1:\newline 9000/9001 & minio-data & 先于 Backend 启动 \\
  ai-service & 项目 Dockerfile & 127.0.0.1:5000 & HTTP health & 先于 Backend 启动 \\
  backend & 项目 Dockerfile & 8080 & HTTP health；启动前迁移 & 等待 MySQL healthy \\
  user-web & Node 构建、Nginx 运行 & 18080 & 无状态静态服务 & Backend 启动后提供访问 \\
  admin-web & Node 构建、Nginx 运行 & 18081 & 无状态静态服务 & Backend 启动后提供访问 \\
\end{longtable}
\normalsize

所有服务采用 \term{restart: unless-stopped}。完整栈配置用于团队联调和课程演示；前端独立
配置用于只构建两个 Web 应用；Podman host 配置用于特定 rootless 网络环境；正式部署叠加
端口收敛配置，取消 MySQL 和 MinIO 的主机端口映射。

\section{部署流程}

\subsection{部署前准备}

部署者确认主机具备 Docker Compose 或 Podman Compose、可用端口和足够磁盘空间，然后从
示例模板生成环境文件。所有空白密码和 token 填写独立随机值，确认前端 API 地址、CORS
来源、MinIO 公共签名地址和 AI 服务地址与目标环境一致。数据库部署还需确认当前
Alembic revision，并在结构变更前制作备份。

\noindent\textbf{完整栈启动命令}\par
\begin{verbatim}
cp deploy/docker/.env.example deploy/docker/.env
# Fill required values in deploy/docker/.env.
docker compose --env-file deploy/docker/.env \
  -f deploy/docker/docker-compose.yml config --quiet
docker compose --env-file deploy/docker/.env \
  -f deploy/docker/docker-compose.yml up --build -d
docker compose --env-file deploy/docker/.env \
  -f deploy/docker/docker-compose.yml ps
\end{verbatim}

先执行 \term{config --quiet} 检查变量展开和 YAML 合并，再构建并后台启动服务。Backend
容器先运行 \term{alembic upgrade head}，迁移成功后才启动 Uvicorn。首次初始化管理员时
显式开启 bootstrap，账号建立并确认可登录后关闭该开关。

\subsection{部署后核验}

\begin{enumerate}
  \item 查看六个服务的状态与启动日志，确认没有重复重启。
  \item 访问 Backend 和 AI Service 的 \path{/health}。
  \item 打开用户端和管理端，检查静态资源及 \path{/api/v1} 代理。
  \item 检查数据库 revision、MinIO bucket 私有策略和对象签名访问。
  \item 使用演示账号执行登录、发布、搜索、匹配、认领、交接和审核的最小走查。
  \item 登记部署时间、commit SHA、镜像标识、数据库 revision、配置版本和执行人。
\end{enumerate}

\subsection{发布变更顺序}

仅修改无状态应用时，先构建固定版本镜像，再逐个替换 AI、Backend 和前端服务；数据库
结构变化时，按照「备份 $\rightarrow$ 迁移 $\rightarrow$ Backend $\rightarrow$ 前端
$\rightarrow$ 走查」的顺序执行。新增字段优先采用兼容方式，使旧应用和新应用在替换期间
都能读取；删除字段先停止应用读写，再在后续版本移除。

\section{运维计划}

\subsection{职责与周期}

\begin{longtable}{L{2.4cm}L{3.1cm}L{5.2cm}L{3.8cm}}
  \caption{运行维护计划}\label{tab:operations-plan}\\
  \toprule
  周期 & 负责人 & 检查内容 & 记录内容 \\
  \midrule
  \endfirsthead
  \multicolumn{4}{c}{\tablename\ \thetable（续）}\\
  \toprule
  周期 & 负责人 & 检查内容 & 记录内容 \\
  \midrule
  \endhead
  \bottomrule
  \endlastfoot
  每次部署 & 发布负责人、后端负责人 & 服务状态、健康接口、revision、主流程走查 & 版本、配置、迁移和核验结果 \\
  每日 & 当值成员 & 服务存活、错误日志、任务状态、磁盘与对象容量 & 异常时间、requestId 和处理动作 \\
  每周 & 运维负责人 & 数据库容量、慢查询、任务积压、日志空间和备份清单 & 趋势、维护项和负责人 \\
  每月 & 项目负责人、运维负责人 & 账号权限、凭据有效期、依赖版本和恢复抽查 & 变更决定和下次检查时间 \\
  重大变更前 & 开发负责人、发布负责人 & 数据备份、迁移演练、版本兼容和回滚入口 & 操作顺序、执行人和确认人 \\
\end{longtable}

\subsection{健康、日志与指标}

MySQL 使用 \term{mysqladmin ping}，Backend 与 AI Service 提供 HTTP health，Compose 根据
健康状态安排 Backend 启动。应用使用 Loguru 输出启动、请求异常、外部服务调用、后台
任务重试和关键状态信息；管理员审核、举报处理和用户管理等操作写入业务操作日志。HTTP
异常响应携带 requestId，便于将客户端反馈与服务日志关联。

\begin{table}[H]
  \centering
  \caption{运行观察项目}
  \begin{tabularx}{\textwidth}{L{3.0cm}L{5.3cm}Y}
    \toprule
    对象 & 观察内容 & 处理方式 \\
    \midrule
    HTTP 服务 & 请求量、延迟、4xx/5xx、健康状态 & 按时间和 requestId 定位接口与版本 \\
    MySQL & 连接数、慢查询、锁等待、存储空间 & 检查查询、索引、连接池和磁盘 \\
    MinIO & 对象容量、上传错误、签名错误 & 核对 bucket、凭据、网络和对象引用 \\
    AI Service & 调用延迟、超时、规则处理比例 & 检查模型配置、服务 token 和调用配额 \\
    后台任务 & PENDING、RUNNING、FAILED 数量与最老任务时间 & 根据任务类型和重试记录处理 \\
    容器主机 & CPU、内存、磁盘和容器重启次数 & 清理日志或镜像并调整资源配置 \\
    \bottomrule
  \end{tabularx}
\end{table}

日志不记录密码、OTP、JWT、AI token、对象存储 secret 或证件原图。运行日志按日期轮转，
演示和联调环境保留最近一周；业务操作日志随数据库备份保存。部署、配置变更和数据恢复
使用独立运维记录，避免与应用日志混在一起。

\subsection{后台任务维护}

匹配、分类和敏感内容处理使用持久任务表。运维人员关注任务状态、尝试次数、下次执行
时间和锁定时间；遗留的 RUNNING 任务按锁超时规则回收，失败任务保留最后错误摘要。
手工重试前先确认输入对象仍存在、业务版本没有变化，并避免同一任务被重复提交。

\section{备份、恢复与回滚}

\subsection{备份计划}

\begin{longtable}{L{2.5cm}L{4.2cm}L{3.4cm}L{4.1cm}}
  \caption{数据与版本备份计划}\label{tab:backup-schedule}\\
  \toprule
  对象 & 方法 & 周期与保留 & 校验方法 \\
  \midrule
  \endfirsthead
  \multicolumn{4}{c}{\tablename\ \thetable（续）}\\
  \toprule
  对象 & 方法 & 周期与保留 & 校验方法 \\
  \midrule
  \endhead
  \bottomrule
  \endlastfoot
  MySQL & 一致性逻辑备份；长期环境同时保留 binlog & 每日，保留 7 日；周备份保留 4 周 & 在隔离数据库中导入并检查 revision \\
  MinIO & 将 bucket 镜像到独立存储并保存对象清单 & 每日增量、每周全量 & 抽样校验 hash，并与数据库对象引用对账 \\
  配置 & 保存无秘密模板、变量名清单和加密凭据副本 & 每次配置变更 & 在新环境中执行配置展开 \\
  镜像与源码 & Git tag、commit SHA、锁文件和不可变镜像 & 每次发布 & 从版本记录取得对应镜像和迁移 \\
  运维记录 & 部署、迁移、备份、恢复和配置变更记录 & 随操作归档，按学期保存 & 检查时间、版本、执行人和核验项 \\
\end{longtable}

备份副本与运行主机分开保存，敏感备份加密并限制访问。每次备份记录开始时间、结束时间、
版本、大小和校验值。课程演示前制作数据库与对象存储快照，演示数据发生修改时可恢复到
统一初始状态。

\subsection{恢复流程}

\begin{enumerate}
  \item 确定目标应用版本、数据库 revision、对象清单和配置版本。
  \item 在隔离环境创建空 MySQL 与空 MinIO，确认备份文件和解密凭据可用。
  \item 恢复数据库和对象，核对数据库中的对象引用是否都能找到对应文件。
  \item 使用与备份匹配的应用镜像启动，确认 revision 后按顺序迁移到目标 head。
  \item 检查健康接口，并走查登录、发布、搜索、认领、交接和管理审核。
  \item 由发布负责人确认后切换服务入口，保存恢复过程和核验结果。
\end{enumerate}

\subsection{应用与数据库回滚}

应用回滚使用上一稳定 commit 对应的固定镜像，不重新构建同名本地镜像。前端静态资源、
Backend 和 AI Service 使用同一发布记录中的兼容版本；切换后重新执行健康检查和最小
业务走查。

数据库优先采用兼容迁移和向前修正。只增加字段的版本可以保持新旧应用兼容；涉及删除
或数据改写时，部署前制作备份并在隔离副本演练。Alembic downgrade 只在确认不会丢失
数据时使用，否则恢复部署前备份，再启动与该 revision 对应的应用版本。

\section{账号与运行安全}

\begin{table}[H]
  \centering
  \caption{运行安全规则}
  \begin{tabularx}{\textwidth}{L{3.0cm}Y Y}
    \toprule
    领域 & 配置规则 & 运维动作 \\
    \midrule
    服务凭据 & JWT、AI token、数据库和 MinIO 使用独立随机值 & 定期更换并确认旧值停止生效 \\
    网络入口 & 数据库、MinIO 和 AI 默认不向公网开放 & 正式环境只通过 HTTPS 网关提供 Web 入口 \\
    对象访问 & bucket 私有，按角色生成短期签名 & 定期检查匿名策略和签名有效期 \\
    管理员账号 & 首次初始化后关闭 bootstrap & 定期检查管理员列表和操作日志 \\
    调试功能 & 共享环境关闭 Debug 和短信调试返回 & 部署核验时检查开关与手机号白名单 \\
    容器镜像 & 按锁文件构建并使用固定版本标识 & 发布时记录镜像标识和源码 commit \\
    \bottomrule
  \end{tabularx}
\end{table}

凭据更换按照「生成新值、更新服务端、更新调用方、检查新值、撤销旧值」的顺序进行。
数据库和 MinIO 为应用创建专用账号，管理员只在维护时使用管理账号。配置和日志的访问
权限按角色分配，课程演示账号不具备修改基础设施配置的权限。

\section{结论}

本文给出了寻迹项目的配置与运维方法：配置通过环境变量和模板分层管理；源代码、依赖、
数据库和镜像使用明确版本标识；持续集成统一执行静态检查、测试与构建；Compose 负责
六服务部署；部署清单、巡检、日志、指标、备份、恢复和回滚构成日常运维流程。

上述规则将「使用哪个版本、如何配置、怎样部署、运行后检查什么」转化为可重复的步骤，
使开发、联调、课程演示和正式部署能够沿用同一套管理方式。
